\section{Attack Framework}
\label{sec:attacks}

We evaluate two attack paradigms: (1) a learned pixel-level CNN preprocessor trained end-to-end, and (2) direct per-video feature-space optimization of a single-frame perturbation.

\subsection{Pixel-Level Learned Preprocessor}
\label{sec:pixel_attack}

We train a residual CNN $g_\theta(V) = V + \delta$ where $\delta = g_\theta(V) - V$ is the perturbation with $\|\delta\|_\infty \leq \varepsilon$. The network is a 75K-parameter ResidualCNN with projection layers to enforce the $L_\infty$ constraint.

The attack loss minimizes:
\begin{align}
  \mathcal{L} = \mathcal{L}_\text{attack} + \lambda_1 \mathcal{L}_\text{LPIPS} + \lambda_2 \mathcal{L}_\text{temporal} + \lambda_3 \mathcal{L}_\text{SSIM}
\end{align}
where $\mathcal{L}_\text{attack}$ is a soft-IoU loss on SAM2's mask prediction with a first-frame centroid prompt, and the constraint terms enforce perceptual quality.

We developed this approach over four rounds of iterative refinement (detailed in \cref{sec:pixel_exps}), progressively fixing: train/eval task mismatch, evaluation bias, SSIM constraint, codec EOT proxy (YUV 4:2:0 + blur + noise + resize), object pointer tokens in memory attention, and temporal positional encoding.

\subsection{Feature-Space Frame-0 Attack}
\label{sec:feature_attack}

Inspired by UAP-SAM2's~\citep{uap_sam2_2025} use of feature-space objectives, we directly optimize a perturbation $\delta_0$ on frame 0 only, targeting SAM2's intermediate representations.

\paragraph{Setup.}
We use projected Adam with 300 steps and 2 independent restarts (best-of-2 kept). The constraint $\|\delta_0\|_\infty \leq \varepsilon$ is enforced by projected gradient. A differentiable H.264 proxy (shared codec state between adversarial and clean branches) is applied during optimization.

\paragraph{Mode B (FPN feature shift).}
The loss maximizes cosine distance between adversarial and clean FPN finest-level features $F \in \RR^{256 \times 64 \times 64}$:
\begin{align}
  J_\text{fpn} = \frac{1}{N_f} \sum_i \left[1 - \frac{F^a_i \cdot F^c_i}{\|F^a_i\|\|F^c_i\|}\right]
\end{align}
This replicates the UAP-SAM2 principle for a direct codec-robustness test.

\paragraph{Mode C+D (maskmem features + object pointer).}
The loss targets the written memory state directly. Let $M_t^a, M_t^c \in \RR^{N_m \times 64}$ be the adversarial and clean maskmem features after codec transform $T$ (shared random state):
\begin{align}
  J_\text{mem\_write}(t) &= \frac{1}{N_m} \sum_i \left[1 - \cos(M_{t,i}^a, M_{t,i}^c)\right] \\
  J_\text{mem\_match}(t) &= \frac{1}{k} \sum_{u=t+1}^{t+k} \left[C(Q_u^c, K_t^c) - C(Q_u^c, K_t^a)\right] \\
  J_\text{ptr}(t) &= (1 - \cos(p_t^a, p_t^c))
\end{align}
where $C(Q, K) = \frac{1}{N_q}\sum_j \tau \log\frac{1}{N_m}\sum_i \exp(q_j^\top k_i / \tau)$ is the log-sum-exp key-query compatibility, and $Q_u^c, K_t^a$ are query tokens from a future clean frame and key tokens from the adversarial memory.

The total loss is:
\begin{align}
  \mathcal{L} = -(\alpha J_\text{mem\_write} + \beta J_\text{mem\_match} + \gamma J_\text{ptr} + \delta J_\text{mask})
\end{align}
with $(\alpha, \beta, \gamma, \delta) = (1.0, 2.0, 0.25, 0.05)$ for Mode C+D and $(1.0, 0.0, 0.25, 0.05)$ for Mode C+D\textsubscript{nm} (no memory matching term).

\paragraph{Evaluation.}
After optimizing $\delta_0$, we inject the adversarial frame 0 into SAM2's \texttt{init\_state()} and propagate with \texttt{propagate\_in\_video()} using a first-frame GT centroid point prompt. The metric $\djfauc$ is computed against the clean-codec baseline (clean frames encoded and decoded with the same H.264 settings).
